\chapter*{Introduction Générale}
\addcontentsline{toc}{chapter}{Introduction Générale}

\par Avec l'augmentation du nombre de véhicules et l'urbanisation croissante, la recherche d'une place de stationnement est devenue un véritable défi dans les grandes villes. Les conducteurs perdent souvent un temps considérable à chercher un parking disponible, tandis que de nombreux propriétaires de parkings privés disposent de places inoccupées qui restent inexploitées. Cette situation engendre une perte de temps, une augmentation de la consommation de carburant ainsi qu'une congestion accrue du trafic urbain.

\par Dans ce contexte, les technologies mobiles et les systèmes intelligents offrent de nouvelles opportunités pour améliorer la gestion du stationnement. En permettant la centralisation des informations, la réservation des places et l'automatisation des paiements, ces solutions contribuent à simplifier l'expérience des utilisateurs tout en optimisant l'exploitation des espaces de stationnement.

\par Ce projet de fin d'études porte sur la conception et le développement de \textbf{TuniPark}, une plateforme intelligente de gestion de parkings composée d'une application mobile destinée aux utilisateurs et d'une interface web réservée aux administrateurs. L'objectif principal du projet est de faciliter la recherche de places de stationnement publiques et privées, de permettre aux propriétaires de proposer leurs propres parkings, de gérer les sessions de stationnement et d'effectuer les paiements de manière sécurisée.

\par La solution développée intègre également un module d'IA permettant de recommander les parkings les plus pertinents en fonction de plusieurs critères, notamment la distance, le taux de disponibilité, les interactions des utilisateurs ainsi que les caractéristiques tarifaires. Cette approche tend à rendre l'expérience utilisateur plus fluide en proposant des recommandations personnalisées.

\par D'un point de vue technique, l'application mobile a été développée avec \textbf{Flutter}, tandis que l'interface d'administration a été réalisée avec \textbf{Angular}. Le Backend repose sur le framework \textbf{NestJS} et utilise \textbf{PostgreSQL} comme système de gestion de base de données via l'ORM \textbf{Prisma}. Les images sont stockées sur \textbf{Cloudinary}, les notifications Push sont assurées par \textbf{Firebase Cloud Messaging (FCM)}, les courriers électroniques sont envoyés à l'aide de \textbf{Nodemailer} et les paiements sont réalisés via la plateforme \textbf{Flouci}. Un microservice d'IA développé en \textbf{Python} avec \textbf{FastAPI} exploite un modèle \textbf{Random Forest} afin de générer les recommandations de parkings.

\par Ce mémoire présente les différentes étapes de réalisation de ce projet. Il commence par une présentation du contexte général et de l'étude de l'existant, puis décrit l'analyse des besoins, l'architecture du système et la conception. Il présente ensuite l'implémentation des différentes fonctionnalités ainsi que les tests réalisés avant de se conclure par une synthèse des résultats obtenus et des perspectives d'amélioration.